\section{Conclusion}
\label{sec:conclusion}

This paper studied the resource requirements of multi-domain rotation rostering from a fresh start.
The rate relaxation gives a lower bound from the maintenance cycles, and the start-up relaxation adds the cost of starting all robots in phase $0$.
The entering theorem constructs a fresh-start schedule that joins a given feasible rotation.
C-RDD uses these results to derive a dispatching rule, construct the start-up and search a short window of deviations.
Its genotype length is independent of the scheduling horizon.

On $160$ test instances, all $800$ C-RDD runs pass the long-run check, and the mean gap at each size ranges from $7.0\%$ to $13.1\%$.
C-RDD has the lowest average rank among the eight methods.
The ablations show that the entering construction and the deviation search both reduce resource use, while the finite-horizon evaluation often accepts schedules that run short later.
These results support using a rotation to organize the long horizon and concentrating evolutionary search on the fresh start.
Future work will study the joint choice of the rotation and its entering schedule, and extend the construction to varying duty and maintenance durations.
